% Generated draft: requires mathematical and source review.
\documentclass[11pt]{article}
\usepackage[margin=25mm]{geometry}
\usepackage{fontspec}
\setmainfont{Latin Modern Roman}
\usepackage{amsmath,amssymb,mathtools,unicode-math}
\setmathfont{Latin Modern Math}
\usepackage{enumitem,needspace,xurl,hyperref}
\setlength{\parindent}{0pt}
\setlength{\parskip}{5pt}
\setlength{\emergencystretch}{3em}
\newcommand{\N}{\mathbb{N}}
\newcommand{\Z}{\mathbb{Z}}
\newcommand{\Q}{\mathbb{Q}}
\newcommand{\R}{\mathbb{R}}
\newcommand{\C}{\mathbb{C}}
\newcommand{\F}{\mathbb{F}}
\newcommand{\A}{\mathbb{A}}
\newcommand{\PP}{\mathbb{P}}
\newcommand{\E}{\mathbb{E}}
\newcommand{\eps}{\varepsilon}
\newcommand{\dd}{\,\mathrm{d}}
\begin{document}
\section*{UnsolvedMath ID: 20001965}
\textbf{Source problem code:} catalogue entry
\subsection{Definitions and mathematical statement}
Conjecture 6.10, attributed to T. Hausel (2003), concerns a hyper-compact hyperk\"ahler manifold \(M\) of real dimension \(4n\). It asks whether the signature of its localized circle-equivariant pairing satisfies
\[
 (-1)^n\sigma(M)\geq0.
\]
The original convention orders coefficients by the sign of their leading term. The full definition of hyper-compactness and of the pairing is not included. Under the aid's convention, localize circle-equivariant cohomology by \(u\), \(\deg u=2\), to coefficients in \(\mathbb Q(u)\), ordered by leading-term sign, and use localization to define integration. A precise choice of ordering and treatment of odd cohomology must still be fixed before identifying every possible meaning of \(\sigma\).

Comment 6.11 says that the conjecture is known for hypertoric manifolds, where it is the Brown--Colbourn inequality for matroid \(h\)-numbers, originally proved in network reliability. Its exact combinatorial formula is not supplied.

The aid's proposed reduction uses the fixed-component hyper-compact interpretation from Hausel's primary paper: a circle-compact holomorphic-symplectic manifold of complex dimension \(2n\), with symplectic form of circle weight one and compact fixed components, under the hypotheses needed for rationalized localization. Let \(\sigma_{\mathrm{ev}}(M)\) be the signature of the symmetric even part of the localized pairing and let \(d_F=\dim_{\mathbb C}F\). The reported formula is
\[
 \sigma_{\mathrm{ev}}(M)=\sum_F(-1)^{n-d_F}\sigma(F),\qquad
 (-1)^n\sigma_{\mathrm{ev}}(M)=\sum_F\sigma(F).
\]
Here \(\sigma(F)\) is the ordinary signature of the compact fixed component; it is zero when its real dimension is \(2\) modulo \(4\), which explains the second equality. Thus this interpretation reduces the conjecture to nonnegativity of the sum of fixed-component signatures, rather than proving that sum nonnegative.

The aid also flags a conflicting finite-critical-point reading, not defined in the extracted conjecture itself. Weight-one symplecticity pairs tangent weights \(a\) and \(1-a\). At an isolated Hamiltonian minimum all weights would have one strict sign, impossible for such integer pairs. With the proper bounded-below moment hypothesis ensuring a minimum, the literal interpretation that all critical points are isolated therefore has no positive-dimensional examples. This does not exclude the fixed-component interpretation.

The record concatenates Section 6.3, ``Abelianization,'' with background credit to S. Martin for the symplectic procedure. Preserve its two reported theorems separately. Let \(A\) be finite dimensional, \(G\) compact, \(T\subset G\) a maximal torus, \(W=N_G(T)/T\), and
\[
 X_G=T^*A\mathbin{/\!/\!/}G,\qquad
 X_T=T^*A\mathbin{/\!/\!/}T.
\]
Reduction levels and smoothness hypotheses are not printed. If both quotients are circle compact, Theorem 6.12 (Hausel--Proudfoot 2003) reports
\[
 \int_{X_G}\widehat\kappa_G(\alpha)
 =\frac1{|W|}\int_{X_T}
       \widehat\kappa_T(\operatorname{res}^G_T\alpha)\,e,
 \qquad \alpha\in\widehat H^*_{S^1\times G}(T^*A),
\]
\[
 e=\prod_{a\in\Delta}a(u-a)
 \in(\operatorname{Sym}\mathfrak t^*)^W\otimes\mathbb Q[u]
 \cong H^*_{S^1\times G}(\mathrm{pt};\mathbb Q),
\]
where \(\Delta\) is the root set and the integrals and hats use the localized convention. Restriction to \(T\) is displayed explicitly rather than treating a \(G\)-class as a \(T\)-class without a map.

Theorem 6.13 (Hausel--Proudfoot 2003) assumes circle compactness, equivariant formality and surjectivity of the intended hyperk\"ahler Kirwan map, and reports the unlocalized ring formula
\[
 H^*_{S^1}(X_G;\mathbb Q)
 \cong H^*_{S^1}(X_T;\mathbb Q)^W/\operatorname{Ann}(e),
\]
where the annihilator is taken in the indicated invariant ring. The source prints the Kirwan source as ordinary \(H^*(T^*A)\); the customary equivariant source and exact form of this hypothesis require checking, and that notation defect is not silently treated as settled.

The source attributes calculation of the toral equivariant ring to Harada--Proudfoot (2002), and describes a hyperpolygon-space completion of the program by Hausel--Proudfoot (2003), involving Konno's spaces and Harada--Proudfoot (2003). Its claim that this was the only completed case is historical context, not a new current classification. These abelianization statements do not prove the main signature inequality.
\subsection{Short English statement}
The signature conjecture is retained with its ordering and definition issues. The aid gives an even-pairing fixed-component formula that reduces it to a sum of ordinary signatures, and explains why an isolated-critical-point interpretation is inconsistent under its hypotheses. The concatenated abelianization integral and ring theorems are preserved separately with their original assumptions and citation labels.
\subsection{Sources}
\begin{itemize}
\item Local source record: catalogue entry, catalogue entry; canonical \texttt{aim-geometry-notes.json}, \texttt{notes[302]}, source item 6.10. Title: ``A fixed-component formula for the hyper-compact signature.''
\item AIM workshop list \emph{Moment maps and surjectivity in various geometries}, supplied reference \url{https://aimath.org/WWN/momentmaps/momentmaps.pdf}. The record includes Conjecture 6.10, Comment 6.11 and attached Section 6.3/Theorems 6.12--6.13. This draft uses the supplied text without accessing a PDF.
\item Supplied attributions: T. Hausel (2003), Brown--Colbourn and network reliability, S. Martin, Hausel--Proudfoot (2003), Harada--Proudfoot (2002), Konno, and Harada--Proudfoot (2003). The aid refers to Hausel's primary fixed-component definition without a full citation. No complete bibliographic entries are supplied.
\item Attempt 1 supplies the even-signature localization reduction and isolated-minimum obstruction. Registry dated 14 August 2026: ``partially\_solved,'' ``PARTIAL-PROGRESS,'' statement status ``exact''; research-run status unknown, attempt partial, rigor ``proved,'' novelty confidence low. The machine-generated, non-peer-reviewed aid requires independent verification before citation.
\end{itemize}
\end{document}
